\begin{proof}

\begin{enumerate}[label={$\langle1\rangle$\arabic*.},leftmargin=*]
\item Choose target coordinates $(v,s)$ in a bidisc about
the section, and let $\Gamma$ be its graph.  Take a
local defining function $\xi$ for $\Gamma$.
On the smooth fourfold the divisor identity is
\[
 g^*\Gamma=m\mathcal D,
\]
where $\mathcal D$ is the prime closure of the
punctured reduced multiple-fibre family.  Indeed
this equality holds away from $s=0$.  The only
possible extra divisorial component would be
$X_0$, but $g_0$ is nonconstant, so $X_0$ is not
mapped into the point $\Gamma(0)$.  The closure
$\mathcal D$ is Cartier because the total space
is smooth.  Thus locally, even if its central
restriction is initially singular or reducible,
\[
 g^*\xi=a^m u
\]
for a local equation $a$ of $\mathcal D$ and a
holomorphic unit $u$.

\item Let $B'\to B$ be the smooth root bidisc
$r^m=\xi$, and let $Y$ be the normalization
of $\mathcal X_B\times_B B'$.  A local $m$th
root of $u$ factors its defining equation as
\[
 r^m-a^m u=
 \prod_{\zeta^m=1}(r-\zeta a u^{1/m}).
\]
Its normalized branches are separate copies of
the original smooth local fourfold.  Therefore
$Y\to\mathcal X_B$ is a finite etale cover,
and $Y$ is smooth.  The deck group permutes these
branches freely, including over $a=0$.
This is an assertion about the normalized cover;
the unnormalized fibre product generally has
intersecting branches.

\item The induced map $g':Y\to B'$ is proper and
equidimensional of relative dimension two.
Since both source and base are smooth, the
two pulled-back base parameters form a regular
sequence in each local ring, giving flatness.
Over $s\ne0$ this is the actual smooth root
torus family.  Its connectedness makes $Y$
connected and its generic fibre connected.
The Stein factor of $g'$ is finite birational
over the normal bidisc and hence is the
bidisc itself, so all fibres are connected.
Equivalently this follows from normality of
$Y$ and the finite birational form of its
Stein factorization.

\item By \ref{push:src:E.2.12} the original central map has no
critical value away from the labelled limits.
Our choice of the section therefore makes
$g'$ smooth away from the origin of $B'$.
It remains to supply a smooth Jacobian over
the full root bidisc.

\item If the normalizing projective transformations
$A_s$ have an invertible limit, the source
coordinate at the finite marked value has
the form
\[
 \xi_{\rm src}=\xi\,u_1(v,s)
\]
with $u_1$ a unit.  Thus its order-$e$ root
is $r^{m/e}u_1^{1/e}$, and the exact finite-root
period model with parameter $c(s)$ extends
smoothly to the full root bidisc.

\item If $A_s$ has rank-one limit, the isolated
section in the hypothesis must be the
preimage of its constant image $t_*$:
all other marked preimages tend to the
exceptional point.  Thus $t_*$ is precisely
the finite marked value under consideration.
Near the isolated section the source
coordinate has divisor equal to that section
plus a positive multiple of $s=0$, so
\[
 \xi_{\rm src}=s^n\xi\,u_1(v,s),\qquad n>0,
\]
with a nowhere-zero unit $u_1$.  Choose a
nearby central regular target point with
$\xi\ne0$.  Its actual smooth torus family
has identity parameter monodromy.  As in
\ref{push:src:E.2.12}, $c(s)$ has no winding and hence
$e\mid n$.  The required holomorphic root
coordinate on $B'$ is now
\[
 s^{n/e}r^{m/e}u_1^{1/e}.
\]
It again gives a smooth reference Jacobian
over $B'$.  On the punctured bidisc its
marking and holomorphic period subbundle
agree with the actual Jacobian by the same
simply-connected argument as \ref{push:src:E.2.12}.

\ref{push:src:E.2.11} makes $g'$ a smooth torsor over the
entire root bidisc.  In particular $r$ is
a submersion coordinate on the central
threefold upstairs.  The unramified deck
cover then shows that the original fibre
has exactly multiplicity $m$.  Restricting
to $r=0$ gives a proper smooth family of
two-tori in the parameter $s$, with free
affine deck action.  Its quotient is a
proper smooth family of the reduced finite
fibres, proving all assertions.

\end{enumerate}
\end{proof}
